\documentclass[11pt,a4paper,leqno]{amsart}%{amsproc}


\usepackage{amssymb}

\usepackage{times}

\usepackage{mathrsfs}

\usepackage{verbatim}


\input{colordvi}


\theoremstyle{plain}

\newtheorem{theorem}{Theorem}[section]

\theoremstyle{remark}

\newtheorem{remark}[theorem]{Remark}

\newtheorem{question}[theorem]{Question}

\newtheorem{example}[theorem]{Example}

\theoremstyle{plain}

\newtheorem{corollary}[theorem]{Corollary}

\newtheorem{lemma}[theorem]{Lemma}

\newtheorem{proposition}[theorem]{Proposition}

\newtheorem{definition}[theorem]{Definition}

\newtheorem{hypothesis}[theorem]{Hypothesis}

\newtheorem{assumption}[theorem]{Assumption}



\newtheorem{Assumption}{Assumption}[section] %%oddzielna numeracja

\renewcommand{\theAssumption}{A.\arabic{Assumption}}



\newenvironment{bew}[2]{\removelastskip\vspace{6pt}\noindent

 {\it Proof  #1.}~\rm#2}{\par\vspace{6pt}}

\newenvironment{sketch}[2]{\removelastskip\vspace{6pt}\noindent

 {\it Main ideas of the proof  #1.}~\rm#2}{\par\vspace{6pt}}





\newenvironment{del}[1]{}{\par\vspace{6pt}}



\numberwithin{equation}{section}





%f\newcommand{\del}[1]{}





\newcommand\Ce{C_{1/2}^\xi(\mathbb{R},\mathrm{E})}

\newcommand\Cx{C_{1/2}^\xi(\mathbb{R},\mathrm{X})}

\newcommand\Oe{\Omega(\xi,\mathrm{E})}

\newcommand\hOe{\hat{\Omega}(\xi,\mathrm{E})}

\newcommand\Ox{\Omega(\xi,\mathrm{X})}

\newcommand\Cox{C_{1/2}(\mathbb{R},\mathrm{X})}



\def\imbed{\hookrightarrow}

\def\todown{\searrow}

\def\toup{\nearrow}

\newcommand{\lb}{\langle}

\newcommand{\rb}{\rangle}



\newcommand{\rf}[1]{(\ref{#1})}

\newcommand\norm[1]{|\!|#1|\!|}

\newcommand\n[1]{|\!|\!|#1|\!|\!|}

\newcommand\SetOf[2]{\bigl\{#1\,\bigm|\,#2\bigr\}}

\newcommand\bo[1]{{\bf #1}}

\newcommand{\embed}{\hookrightarrow}

\newcommand\lin{\operatorname{lin}}

\newcommand\PG{\operatorname{PG}}

\newcommand\sgn{\operatorname{sgn}}

\newcommand\grad{\operatorname{grad}}

\newcommand\diver{\operatorname{div}}

\newcommand\curl{\operatorname{curl}}

\newcommand\rot{\operatorname{rot}}

\newcommand\supp{\operatorname{supp}}

\newcommand\rank{\operatorname{rank}}

\newcommand\Res{\operatorname{Res}}

\newcommand\AG{\operatorname{AG}}

\newcommand\id{\operatorname{id}}

\newcommand\vnl{(\!(}

\newcommand\vnr{)\!)}

\newcommand\Vertt{|\!|}

\newcommand\den{:\!=}

\newcommand{\ff}{\emph{focus-focus}}

\newcommand{\U}{\mathcal{U}}

\newcommand{\M}{\mathcal{M}}

\newcommand{\eps}{\varepsilon}

\newcommand{\al}{\alpha}

\newcommand{\la}{\lambda}

\def\Nat{{\mathbb N}}

\def\Qnu{{\mathbb Q}}

\def\Rnu{{\mathbb R}}

\def\Cnu{{\mathbb C}}

\def\Rn{{\Rnu^n~}}

\def\Rt{{\Rnu^1~}}

\def\Rd{{\Rnu^d~}}

\def\Rk{{\Rnu^k~}}

\def\Rnus{\mbox{\footnotesize\bf R}}

\def\Re{\mbox{Re\,}}

\def\Im{\mbox{Im\,}}

\def\im{\mbox{im\,}}

\def\liml{\lim\limits}

\def\suml{\sum\limits}

\def\supl{\sup\limits}

\def\maxl{\max\limits}

\def\minl{\min\limits}

\def\cupl{\mathop{\cup}\limits}

\def\capl{\mathop{\cap}\limits}

\def\intl{\int\limits}

\def\prodl{\prod\limits}

\def\supl{\mathop{\sup}\limits}



\newcommand{\bmu}{\mbox{\boldmath{$\mu$}}}



\newcommand{\cit}[1]{\cite{[#1]}}

\sloppy


\begin{document}
\title{Global evolution of random vortex filament equation}

\author{Z. Brze\'{z}niak$^\dag$, M. Gubinelli$^\ddag$ and
M. Neklyudov$^\dag$}

\address{
$^\dag$Department of Mathematics, University of York, Heslington,
UK
\\
$^\ddag$CEREMADE, Paris, France
}
\date{}


\date{\today}

\begin{abstract}
In this article we prove existence of global solution for random
vortex filament equation. Our work gives a positive answer to a question left open in recent publications: Berselli and  Gubinelli \cite{[Bers+Gub_2007]}  showed the existence of global solution for a smooth  initial condition  while  Bessaih, Gubinelli, Russo \cite{[Bess+Gub_2005]} proved the existence of a local solution for a general  initial condition.
\end{abstract}


\maketitle

In this article we prove the existence of a global solution for the following random
vortex filament equation
\begin{eqnarray}
\frac{d\gamma}{dt} &=& u^{\gamma(t)}(\gamma(t)), \;\; t\in [0,T] \label{eqn:RandomFilam-1}\\
\gamma(0) &=& \gamma_0,\label{eqn:RandomFilam-2}
\end{eqnarray}
where the initial condition $\gamma_0$ is a  geometric $\nu$-rough path (for some  $\nu\in(\frac{1}{3},1)$),  see Assumption \ref{ass:RoughPathApproximability}. Here $\gamma:[0,T]\to \mathcal{D}_{\gamma_0}\subset\mathcal{C}$ is some trajectory in
the subset $\mathcal{D}_{\gamma_0}$ of $\mathcal{C}$ of continuous closed curves in
$\Rnu^3$, $u^Y$, $Y\in \mathcal{D}_{\gamma_0}\subset\mathcal{C}$ is a vector field
given by
\begin{equation}
u^Y(x)=\intl_Y\nabla\phi(x-y)\times dy.\label{eqn:RandomFilam-3}
\end{equation}
where $\phi:\Rnu^3\to\Rnu$ is a smooth function which satisfies
certain assumptions (see Hypothesis \ref{hyp:PhiAssumptions}).
Exact meaning of the line integral above and set $\mathcal{D}_{\gamma_0}$ we consider
will be explained below. This equation appear in fluid dynamics in
the theory of three dimensional Euler equations.
It is well known that for two dimensional Euler equations vorticity $\vec{\omega}=\curl \vec{u}$ is transported along the flow of the liquid.
The situation changes drastically in three dimensional case. Additional "stretching" term in the equation defining vorticity
leads to possibility of blow up of the vorticity.
Furthermore, result of Beale,
Kato, and Majda \cite{[BKM-1984]} suggest that possible
singularity of Euler equations appear when vorticity field of
fluid blows up. Consequently, understanding of the behaviour of vorticity of ideal fluid is
one of the most important problems in fluid dynamics.

The properties of the motion of the vorticity has been studied for the last 150 years starting from the works of
Helmholtz \cite{Helm1858} and Kelvin \cite{Kelvin1883}. It has been suggested by Kelvin %(not sure???) 
to use Biot-Savart law
\[
\vec{u}(x)=\int\frac{\vec{x}-\vec{y}}{\vert \vec{x}-\vec{y}\vert^3}\times\vec{\omega}(\vec{y})d\vec{y},x,y\in\mathbb{R}^3
\]
combined with assumption that vorticity is supported by some smooth curve $\gamma$
\[
\omega(\vec{x},t)=\Gamma\int\limits_0^1\delta(\vec{x}-\gamma(s,t))\frac{\partial \gamma(s,t)}{\partial s}ds,x\in\mathbb{R}^3,t\geq 0
\]
and definition of the flow
\begin{equation}\label{equa21}
\left\{
\begin{array}{rcl}
\frac{d \vec{X}_t(\vec{x})}{dt}&=&\vec{u}(\vec{X}_t(\vec{x}),t),\;\;t\geq 0,\\
\vec{X}_0(\vec{x})&=&\vec{x}.
\end{array}
\right.
\end{equation}
to formally deduce random filament equation
\begin{equation}\label{eqn:RandomFilamClassic}
\frac{\partial \gamma}{\partial t}(s,t)=-\frac{\Gamma}{4\pi}\int\limits_0^1\frac{\gamma(s,t)-\gamma(r,t)}{\vert \gamma(s,t)-\gamma(r,t)\vert^3}\times\frac{\partial\gamma(r,t)}{\partial r}dr.
\end{equation}
Assumption that vorticity is supported by some curve is substantiated by some numerical simulations of 3D turbulent fluids
which show that regions where vorticity is big have a form of
"filament", see, for instance \cite{[Bell+Markus_1992]},
\cite{[VincMeneguzzi-1991]}.

The equation \eqref{eqn:RandomFilamClassic} has singularity when $r$ is close to $s$ and initial curve $\gamma$ is smooth.
As a consequence, the energy of the solution of this equation given by formula
\[
E(t)=\frac{\Gamma^2}{8\pi}\int\limits_0^1\int\limits_0^1\frac{1}{\vert \gamma(s,t)-\gamma(r,t)\vert }\frac{\partial \gamma(r,t)}{\partial r}\cdot\frac{\partial \gamma(s,t)}{\partial s}dsdr,t\geq 0
\]
is infinite for any smooth curve $\gamma(\cdot,t)$.
Hence different methods are employed to avoid singularity. For instance, Gallavotti \cite{Gallavotti} motivated by finiteness of the energy integral
for Brownian Motion\footnote{defined as double stochastic integral},
considers non smooth initial curves $\gamma_0$.
Rosenhead \cite{[Rosenhead-1930]}  has suggested to use the following model
\begin{equation}\label{eqn:RandomFilamRosenhead}
\frac{\partial \gamma}{\partial t}(s,t)=-\frac{\Gamma}{4\pi}\int\limits_0^1\frac{\gamma(s,t)-\gamma(r,t)}{(\vert \gamma(s,t)-\gamma(r,t)\vert^2+\mu^2)^{3/2}}\times\frac{\partial\gamma(r,t)}{\partial r}dr.
\end{equation}
The equation (\ref{eqn:RandomFilam-1}-\eqref{eqn:RandomFilam-3}) has been introduced  by Berselli, Gubinelli \cite{[Bers+Gub_2007]} (see also \cite{[Bers+Bess_2002]}).
It contains equation \eqref{eqn:RandomFilamRosenhead} considered by Rosenhead as a very particular case when
\[
\phi(\vec{x})=\frac{\Gamma}{(\vert \vec{x}\vert^2+\mu^2)^{1/2}},\vec{x}\in\mathbb{R}^3,\mu>0.
\]
Berselli and  Gubinelli \cite{[Bers+Gub_2007]}  showed the existence of global solution to  equation (\ref{eqn:RandomFilam-1}-\ref{eqn:RandomFilam-3}) if initial condition is a smooth curve. The existence of a local solution to  equation (\ref{eqn:RandomFilam-1}-\ref{eqn:RandomFilam-3}) when initial condition is a non smooth curve has been established in Bessaih, Gubinelli, Russo \cite{[Bess+Gub_2005]}. Our aim is to show global existence of solution of the equation (\ref{eqn:RandomFilam-1}-\ref{eqn:RandomFilam-3}) when initial
condition is a non smooth curve.

We will work in the framework of rough path
theory developed by T. J. Lyons and co-authours, see
\cite{[Lyons-1998]}, \cite{[LyonsQian-2002]} and references
therein, and assume that initial condition is a closed curve of
H\"{o}lder class with exponent $\nu\in (\frac{1}{3},1]$.

\section{Notation}
In this section we present framework of Gubinelli, see
\cite{[Gubinelli-2004]}.
Let $V$ be a fixed Banach space. We define the following two objects:
\begin{eqnarray*}
C_n(V)&=&\{f\in C((S^1)^n,V):\; f(t,t,\ldots,t)=0,t\in
S^1\},\\
C_*(V)&=&\cupl_{k\in\Nat}C_k(V).
\end{eqnarray*}
Operator
\begin{eqnarray*}
\delta&:&C_n(V)\to C_{n+1}(V),\\
(\delta g)(t_1,\ldots,t_{n+1})&=&\suml_{i=1}^{n+1}(-1)^ig(t_1,\ldots\hat{t_i}\ldots,t_{n+1}),n\in\Nat
\end{eqnarray*}
satisfies the following fundamental property
$$
\delta\delta=0,
$$
where $\delta\delta$ is understood as an operator from $C_n(V)$ to
$C_{n+2}(V)$. Thus $\delta$ induces a complex and we can denote
\begin{eqnarray*}
\mathcal{Z}C_k(V)&=&C_k(V)\cap\ker\delta,\\
\mathcal{B}C_k(V)&=&C_k(V)\cap\im\delta.
\end{eqnarray*}
To avoid confusion we will use notation $\delta_n$ for operator
$\delta:C_n(V)\to C_{n+1}(V)$. Furthermore, it can be noticed that
$$
\im\delta_n=\ker\delta_{n+1},
$$
i.e. $\mathcal{Z}C_{k+1}(V)=\mathcal{B}C_k(V)$.

We will  mainly consider the cases $n=1$ and $n=2$. Then operator $\delta$
has following form
$$
\delta_1g(t,s)=g(t)-g(s),\delta_2h(t,u,s)=h(t,s)-h(t,u)-h(u,s),t,u,s\in
S^1.
$$
and we will use special topology in spaces $C_2(V)$ and $C_3(V)$.
Let
\begin{eqnarray*}
\vert f\vert_{\mu}&=&\supl_{a,b\in S^1}\frac{\vert f(a,b)\vert_{V}}{\vert a-b\vert^{\mu}},\;\;f\in
C_2(V),
\\
C_2^{\mu}(V)&=&\{f\in C_2(V):\, \vert f\vert_{\mu}<\infty\}
\\
\vert g\vert_{\rho,\mu}&=&\supl_{a,b,c\in
S^1}\frac{\vert g(a,b,c)\vert_{V}}{\vert a-b\vert^{\rho}\vert b-c\vert^{\mu}},\;\;g\in
C_3(V),
\\
C_3^{\rho,\mu}(V)&=&\{f\in C_3(V):\, \vert f\vert_{\rho,\mu}<\infty\},
\\
C_3^{\mu}(V)&=&\{f\in
C_3(V): f=\suml_ih_i,\vert f\vert_{\mu}=\suml_i\vert h_i\vert_{\rho_i,\mu-\rho_i}<\infty,\rho_i\in
(0,\mu)\},
\\
\mathcal{Z}C_3^{\mu}(V)&=&C_3^{\mu}(V)\cap\mathcal{Z}C_3(V),
\\
\mathcal{Z}C_3^{1+}(V)&=&\cupl_{\mu>1}\mathcal{Z}C_3^{\mu}(V),\\
C_2^{1+}(V)&=&\cupl_{\mu>1}C_2^{\mu}(V).
\end{eqnarray*}

Then following fundamental proposition has been proved in
\cite{[Gubinelli-2004]}:
\begin{proposition}
There exists an unique linear map
$\Lambda:\mathcal{Z}C_3^{1+}(V)\to C_2^{1+}(V)$ such that
$$
\delta\Lambda=\id_{\mathcal{Z}C_3^{1+}(V)}.
$$
Furthermore, for any $\mu>1$, this map is continuous from
$\mathcal{Z}C_3^{\mu}(V)$ to $C_2^{\mu}(V)$ and we have
$$
\Vert \Lambda h\Vert_{\mu}\leq
\frac{1}{2^{\mu}-2}\Vert h\Vert_{\mu},h\in\mathcal{Z}C_3^{1+}(V).
$$
\end{proposition}
Now, we define class of paths for which rough path integral will
be defined.
\begin{definition}
Let us fix $X\in C^{\nu}(S^1,V)$. %, $X\notin
%C^{2\nu}(S^1,\Rnu^3)$.
We say that path $Y\in C(S^1,V)$ is weakly controlled by $X$ if
there exist functions $Z\in C^{\nu}(S^1,L(V,V))$ and $R\in
C_2^{2\nu}(V)$ such that
\begin{equation}
Y(\xi)-Y(\eta) =Z(\eta)(X(\xi)-X(\eta))+R(\xi,\eta),\xi,\eta\in
S^1,\label{eqn:WeakControlDef-1}
\end{equation}
\end{definition}
Let $\mathcal{D}_X$ be the set of pairs $(Y,Z)$, where $Y\in
C(S^1,V)$ is a path weakly controlled by $X$, and $Z\in
C^{\nu}(S^1,L(V,V))$ is such that $R\in C_2^{2\nu}(V)$, where $R$
is defined by representation \eqref{eqn:WeakControlDef-1}. We
notice that this is vector space. Define the semi-norm in
$\mathcal{D}_X$ as follows
\begin{equation}
\Vert (Y,Z)\Vert_{\mathcal{D}_X}=\vert Z\vert_{C^{\nu}}+\Vert R\Vert_{C_2^{2\nu}},\label{eqn:D_XNorm-1}
\end{equation}
where
\begin{equation}
R(\xi,\eta)=Y(\xi)-Y(\eta)-Z(\eta)(X(\xi)-X(\eta)),\xi,\eta\in
S^1.\label{eqn:WeakControlDef-1'}
\end{equation}
Furthermore, define the following norm
\begin{equation}
\Vert (Y,Z)\Vert_{\mathcal{D}_X}^*=\Vert Y\Vert_{\mathcal{D}_X}+\vert Y\vert_{C(S^1,V)}.\label{eqn:D_XNorm-2}
\end{equation}
Then, one can prove that
$(\mathcal{D}_X,\Vert \cdot\Vert_{\mathcal{D}_X}^*)$ is a Banach space.
From now on we will denote elements of $\mathcal{D}_X$ by $(Y,Y')$
and the corresponding $R$ will be denoted by $R^Y$. We will often
omit to specify $Y'$ when it is clear from the context and write
$\Vert Y\Vert_{\mathcal{D}_X}$ instead of $\Vert (Y,Y')\Vert_{\mathcal{D}_X}$.

\section{Definition and Properties of Rough Path integrals}
In this section we define rough path integral and state some of
its properties. We mainly follow \cite{[Gubinelli-2004]} and
\cite{[Bess+Gub_2005]}.
%\begin{example}
%insert Brownian loops
%\end{example}


\begin{definition}
Let $\Pi:\mathcal{D}_X\ni (Y,Z)\mapsto Y\in C(S^1,\Rnu^3)$ be the
natural projection.
\end{definition}
We will need following properties of $\mathcal{D}_X$, see
\cite{[Gubinelli-2004]}.
\begin{lemma}\label{lem:HolderProp}
$\Pi(\mathcal{D}_X)\subset C^{\nu}(S^1,\Rnu^3)$.
\end{lemma}
\begin{proof}[ Proof of Lemma \ref{lem:HolderProp}]
Immediately follows from inequality
\begin{equation}
\Vert Y\Vert_{C^{\nu}}\leq
\Vert Y\Vert_{\mathcal{D}_X}^*(1+\Vert X\Vert_{C^{\nu}}).\label{eqn:HolderIneq}
\end{equation}
\end{proof}
\begin{lemma}\label{lem:RegularMap}
Let $\phi\in C^2(\Rnu^3,\Rnu^3)$ and $(Y,Z)\in \mathcal{D}_X$.
Then
\begin{equation}
(W,W'):=(\phi(Y),\phi'(Y)Z)\in
\mathcal{D}_X\label{eqn:RegularMap-1}
\end{equation}
and the remainder has the following representation
\begin{eqnarray}
R^{W}(\xi,\eta)=\phi'(Y(\xi))R(\xi,\eta)+
(Y(\eta)-Y(\xi))\nonumber
\end{eqnarray}
\begin{eqnarray}
+\intl_0^1[\nabla\phi(Y(\xi)+r(Y(\eta)-Y(\xi)))-\nabla\phi(Y(\xi))]dr,\xi,\eta\in
S^1.\label{eqn:RegularMap-2}
\end{eqnarray}
where $R$ is the remainder for $Y$ w.r.t. $X$ given by
\eqref{eqn:WeakControlDef-1'}. Furthermore, there exists a
constant $K\geq 1$ such that
\begin{equation}
\Vert \phi(Y)\Vert_{\mathcal{D}_X}\leq
K\Vert \nabla\phi\Vert_{C^1}\Vert Y\Vert_{\mathcal{D}_X}(1+\Vert Y\Vert_{\mathcal{D}_X})(1+\Vert X\Vert_{C^{\nu}})^2.\label{eqn:RegularMap-3}
\end{equation}
Moreover, if $(\tilde{Y},\tilde{Z})\in \mathcal{D}_{\tilde{X}}$
and
$$
(\tilde{W},\tilde{W}'):=(\phi(\tilde{Y}),\phi'(\tilde{Y})\tilde{Z})
$$
then
\begin{eqnarray}
\vert W'-\tilde{W}'\vert_{C^{\nu}}+\vert R^W-R^{\tilde{W}}\vert_{C_2^{2\nu}}+\vert W-\tilde{W}\vert_{C^{\nu}}\leq\nonumber\\
C(\vert X-\tilde{X}\vert_{C^{\nu}}+\vert Y'-\tilde{Y}'\vert_{C^{\nu}}+\vert R^Y-R^{\tilde{Y}}\vert_{C_2^{2\nu}}+\vert Y-\tilde{Y}\vert_{C^{\nu}})\label{eqn:RegularMap-4}
\end{eqnarray}
with
\begin{equation}
C=K\Vert \phi\Vert_{C^3}(1+\Vert X\Vert_{C^{\nu}}+\Vert \tilde{X}\Vert_{C^{\nu}})^3\vert (1+\Vert Y\Vert_{\mathcal{D}_X}+\Vert \tilde{Y}\Vert_{\mathcal{D}_{\tilde{X}}})^2.\label{eqn:RegularMap-4'}
\end{equation}
In the case $X=\tilde{X}$ we have
\begin{eqnarray}
\Vert \phi(Y)-\phi(\tilde{Y})\Vert_{\mathcal{D}_X}\leq
K\Vert \nabla\phi\Vert_{C^2}\Vert Y\Vert_{\mathcal{D}_X}\nonumber\\
(1+\Vert Y\Vert_{\mathcal{D}_X}+\Vert \tilde{Y}\Vert_{\mathcal{D}_X})^2(1+\Vert X\Vert_{C^{\nu}})^4
\Vert Y-\tilde{Y}\Vert_{\mathcal{D}_X}.\label{eqn:RegularMap-4''}
\end{eqnarray}
\end{lemma}
\begin{proof}[ Proof of Lemma \ref{lem:RegularMap}]
See \cite{[Gubinelli-2004]}, Proposition $4$ for all statements of
the Lemma, except \eqref{eqn:RegularMap-2} (which is actually also
proven, though not stated explicitly). Let us show
\eqref{eqn:RegularMap-2}. Denote
$y(r)=Y(\xi)+r(Y(\eta)-Y(\xi)),r\in[0,1]$. Then
\begin{eqnarray}
\label{eqn:RegularMap-2proof}
&&
\phi(y(1))-\phi(y(0))=\intl_0^1\phi'(y(r))y'(r)dr
\\
\nonumber
&=&\suml_k(Y^k(\eta)-Y^k(\xi))\intl_0^1\frac{\partial\phi}{\partial
x_k}(y(r))dr\\
&=&\suml_k\frac{\partial\phi}{\partial
x_k}(Y(\xi))(Y^k(\eta)-Y^k(\xi))\nonumber
\\
&+&\suml_k(Y^k(\eta)-Y^k(\xi))\intl_0^1[\frac{\partial\phi}{\partial
x_k}(y(r))-\frac{\partial\phi}{\partial x_k}(Y(\xi))]dr\nonumber
\\
&=&\suml_{k,l}\frac{\partial\phi}{\partial
x_k}(Y(\xi))(Y')^{kl}(X^l(\eta)-X^l(\xi))
+\suml_{k}\frac{\partial\phi}{\partial
x_k}(Y(\xi))(R^Y)^k(\xi,\eta)\nonumber
\\
&+&\suml_k(Y^k(\eta)-Y^k(\xi))\intl_0^1[\frac{\partial\phi}{\partial
x_k}(y(r))-\frac{\partial\phi}{\partial
x_k}(Y(\xi))]dr,
\nonumber
\end{eqnarray}
and the result follows.
\end{proof}
Now we define integral of path $Y$ weakly controlled by $X$ w.r.t.
another path $Z$, weakly controlled by $X$. We will need one more
definition.
\begin{definition}\label{def:RoughPath}
Let $\nu>\frac{1}{3}$. We say that couple
$\mathbb{X}=(X,\mathbb{X}^2)$, $X\in C^{\nu}(S^1,\Rnu^3)$,
$\mathbb{X}^2\in C_2^{2\nu}(L(\Rnu^3,\Rnu^3))$ is a $\nu$-rough
path if the following condition is satisfied:
\begin{equation}\label{eqn:CompatCondition}
\mathbb{X}^{2}(\xi,\rho)-\mathbb{X}^{2}(\xi,\eta)-\mathbb{X}^{2}(\eta,\rho)
=(X(\xi)-X(\eta))\otimes(X(\eta)-X(\rho)), \xi,\eta,\rho\in S^1
\end{equation}
\end{definition}
\begin{remark}
If $\nu>1$ and $\mathbb{X}$ is a $\nu$-rough path, then
$\mathbb{X}$ is identically pair of constants $(X(0),0)$. Indeed,
in this case $X$ is H\"older function with exponent more than $1$
i.e. constant $X(0)$. Hence, $\mathbb{X}^{2}=0$.
\end{remark}
\begin{remark}
If $\nu\in(\frac{1}{2},1]$ then $\mathbb{X}^2$, the second
component of a $\nu$-rough path $\mathbb{X}=(X,\mathbb{X}^2)$, is
uniquely determined by its first component. Indeed,
\begin{equation}
\mathbb{X}^{2,ij}(\xi,\eta)=\intl_{\xi}^{\eta}(X^i_{\rho}-X^i_{\eta})dX^j_{\rho},\xi,\eta\in
S^1,i,j=1,2,3\label{eqn:YoungIntegralCase}
\end{equation}
where the integral is understood in the sense of Young, see
\cite{[Young-1936]}. One can show that $\mathbb{X}^{2}$ defined by
formula \eqref{eqn:YoungIntegralCase} satisfies conditions of
Definition \ref{def:RoughPath}. Let us show uniqueness of
$\mathbb{X}^{2}$. Assume that there exists another
$\mathbb{X}_1^{2}$ which satisfies definition \ref{def:RoughPath}.
Put $G(\xi)=\mathbb{X}^{2}(\xi,0)-\mathbb{X}_1^{2}(\xi,0)$. Then
by condition \ref{eqn:CompatCondition}
$$
\mathbb{X}^{2}(\xi,\rho)-\mathbb{X}_1^{2}(\xi,\rho)=G(\xi)-G(\rho),
$$
and, since $\mathbb{X}^{2}\in C_2^{2\nu}$, $G$ is a H\"{o}lder
function of order bigger  than $1$. Hence,  $G=0$. Therefore,
$\mathbb{X}_1^{2}=\mathbb{X}^{2}$.
\end{remark}
Note that by identity \eqref{eqn:CompatCondition} it follows that
$\mathbb{X}^2(\xi,\xi)=0,\xi\in S^1$.
\begin{assumption}\label{ass:RoughPathApproximability}
We say that our $\nu$-rough path $(X,\mathbb{X}^2)$  is an
geometric $\nu$-rough path if there exist a sequence
$(X_n,\mathbb{X}_n^2)$ such that
$$
X_n\in C^{\infty}(S^1,\Rnu^3),
$$
$$
\mathbb{X}_n^{2,ij}(\xi,\eta)=\intl_{\xi}^{\eta}(X_n^i({\rho})-X_n^i({\eta}))dX_n^j({\rho}),\xi,\eta\in
S^1,i,j=1,2,3,
$$
and
\begin{equation}
\lim_{n\to\infty}\big[\vert X_n-X\vert_{C^{\nu}}+\vert \mathbb{X}_n^2-\mathbb{X}^2\vert_{C_2^{2\nu}}\big]=
0.\label{eqn:RoughPathTopology}
\end{equation}
\end{assumption}
\begin{example}
Let $\{B_t\}_{t\in [0,1]}$ be standard three dimensional Brownian
bridge such that $B_0=B_1=x_0$ and let $\mathbb{B}^{2,ij}$ be the
area process
$$
\mathbb{B}^{2,ij}(\xi,\eta)=\intl_{\xi}^{\eta}(B^i_{\rho}-B^i_{\eta})dB^j_{\rho},i,j=1,2,3
$$
where the integral can be understood either in the Stratonovich or in the It\^o
sense. Then, this couple $(B,\mathbb{B}^2)$ is a $\nu$-rough path see \cite[p.1849]{[Bess+Gub_2005]}. Moreover, if the integral is understood in Stratonovich sense it is
geometric $\nu$-rough path.  Indeed, it follows from Theorem
3.1 in \cite{[Friz-2005]} that we can approximate $X$ with
piecewise linear dyadic $X_n'$ in the sense of assumption
\ref{ass:RoughPathApproximability}a.s..
\end{example}
From now on we suppose that the geometric $\nu$-rough path
$\mathbb{X}=(X,\mathbb{X}^2)$ and the corresponding Banach space
$\mathcal{D}_X$ are fixed.
\begin{lemma}\label{lem:IntegralDef}
Let $\pi=\{\xi_0=\xi<\xi_1<\ldots<\xi_n=\eta\}$ be a finite
partition of $[\xi,\eta]$ and $d(\pi)=\supl_i\vert \xi_{i+1}-\xi_i\vert $ is
a mesh of $\pi$. If $Y,Z\in \mathcal{D}_X$ then the limit
\begin{equation}
\liml_{d(\pi)\to
0}\sum_{i=0}^{n-1}[Y(\xi_i)(Z(\xi_{i+1})-Z(\xi_i))+Y'(\xi_i)Z'(\xi_i)\mathbb{X}^2(\xi_{i+1},\xi_i)]\label{eqn:IntegralDef-1}
\end{equation}
exists and is denoted by definition by
$$
\intl_{\xi}^{\eta}YdZ.
$$
\end{lemma}
\begin{proof}[ Proof of Lemma \ref{lem:IntegralDef}]
See \cite{[Gubinelli-2004]}, Theorem $1$.
\end{proof}
\begin{remark}\label{rem:SmoothCase}
In the case of $\nu>\frac{1}{2}$ the line integral defined in the
Lemma \ref{lem:IntegralDef} is reduced to the Young definition of the
line integral $\int YdZ$. Indeed, it is enough to notice that
second term in formula \eqref{eqn:IntegralDef-1} is of the order
$O(\vert \xi_{i+1}-\xi_i\vert^{2\nu})$, $2\nu>1$. Obviously, line integral
does not depend upon $Y'$, $Z'$ in this case.
\end{remark}
\begin{lemma}\label{lem:IntegralProp}
Assume $Y,W\in \mathcal{D}_X$, $\tilde{Y},\tilde{W}\in
\mathcal{D}_{\tilde{X}}$. Define maps $Q,\tilde{Q}:(S^1)^2\to\Rnu$ by
the following identities
\begin{equation}
Q(\eta,\xi):=\intl_{\xi}^{\eta}YdW-Y(\xi)(W(\eta)-W(\xi))-Y'(\xi)W'(\xi)\mathbb{X}^2(\eta,\xi),\eta,\xi\in
S^1,\label{eqn:IntegralDef-2}
\end{equation}
\begin{equation}
\tilde{Q}(\eta,\xi):=\intl_{\xi}^{\eta}\tilde{Y}d\tilde{W}-\tilde{Y}(\xi)(\tilde{W}(\eta)-\tilde{W}(\xi))-\tilde{Y}'(\xi)\tilde{W}'(\xi)\tilde{\mathbb{X}}^2(\eta,\xi),\eta,\xi\in
S^1. \label{eqn:IntegralDef-4}
\end{equation}
Then $Q,\tilde{Q}\in C_2^{3\nu}$. \\
Moreover, there exists constant
$C=C(\nu)>0$ such that for all $Y,W\in \mathcal{D}_X$
\begin{equation}
\Vert Q\Vert_{C_2^{3\nu}}\leq
C(1+\Vert X\Vert_{C^{\nu}}+\Vert \mathbb{X}^2\Vert_{C_2^{2\nu}})\Vert Y\Vert_{\mathcal{D}_X}\Vert W\Vert_{\mathcal{D}_X}.\label{eqn:IntegralDef-3}
\end{equation}
Furthermore,
\begin{multline}\label{eqn:IntegralDef-5}
\Vert Q-\tilde{Q}\Vert_{C_2^{3\nu}}\leq
C(1+\Vert X\Vert_{C^{\nu}}+\Vert \mathbb{X}^2\Vert_{C_2^{2\nu}})\\
((\Vert Y\Vert_{\mathcal{D}_X}+\Vert \tilde{Y}\Vert_{\mathcal{D}_{\tilde{X}}})\eps_W+
(\Vert W\Vert_{\mathcal{D}_X}+\Vert \tilde{W}\Vert_{\mathcal{D}_{\tilde{X}}})\eps_Y+\eps_X).
\end{multline}
where
\[
\eps_Y=\vert Y'-\tilde{Y}'\vert_{C^{\nu}}+\vert R^Y-R^{\tilde{Y}}\vert_{C_2^{2\nu}}+\vert Y-\tilde{Y}\vert_{C^{\nu}},
\]
\[
\eps_W=\vert W'-\tilde{W}'\vert_{C^{\nu}}+\vert R^W-R^{\tilde{W}}\vert_{C_2^{2\nu}}+\vert W-\tilde{W}\vert_{C^{\nu}},
\]
\[
\eps_X=(\Vert Y\Vert_{\mathcal{D}_X}+\Vert \tilde{Y}\Vert_{\mathcal{D}_{\tilde{X}}})(\Vert W\Vert_{\mathcal{D}_X}+\Vert \tilde{W}\Vert_{\mathcal{D}_{\tilde{X}}})
(\vert X-\tilde{X}\vert_{C^{\nu}}+\vert \mathbb{X}^2-\mathbb{\tilde{X}}^2\vert_{C_2^{2\nu}}).
\]
\end{lemma}
\begin{proof}[ Proof of Lemma \ref{lem:IntegralProp}]
See \cite{[Gubinelli-2004]}, Theorem 1. For formula
\eqref{eqn:IntegralDef-5} see   \cite{[Gubinelli-2004]}, p.104,
formula (27).
\end{proof}
By Lemmata \ref{lem:IntegralDef} and \ref{lem:RegularMap} for any
$A\in C^2(\Rnu^3,L(\Rnu^3,\Rnu^3)),Y\in\mathcal{D}_X$ we can a
define a map  $V^Y:\Rnu^3\to\Rnu$  by invoking rough path integral as follows
\begin{equation}
V^Y(x):=\intl_{S^1}A(x-Y)dY,x\in\Rnu^3.
\end{equation}
We have following bounds on its regularity:
\begin{lemma}\label{lem:ParametricRoughPath}
Let $Y\in \mathcal{D}_X$, $\tilde{Y}\in \mathcal{D}_{\tilde{X}}$,
then there exists $C_1=C_1(\nu),C_2=C_2(\mathbb{X})$ such that for
any integer $n\geq 0$,
\begin{equation}
\Vert \nabla^nV^Y\Vert_{L^{\infty}}\leq
4C_1C_2^3\Vert \nabla^{n+1}A\Vert_{C^1}\Vert Y\Vert_{\mathcal{D}_X}^2(1+\Vert Y\Vert_{\mathcal{D}_X})\label{eqn:ParametricRoughPath-1}
\end{equation}
and
\begin{multline}\label{eqn:ParametricRoughPath-2general}
\Vert \nabla^nV^Y-\nabla^nV^{\tilde{Y}}\Vert_{L^{\infty}}\leq
C(\nu)\vert A\vert_{C^{n+3}}C_X^4(1+\Vert Y\Vert_{\mathcal{D}_X}+\Vert \tilde{Y}\Vert_{\mathcal{D}_{\tilde{X}}})^3\\
(\vert X-\tilde{X}\vert_{C^{\nu}}+\vert \mathbb{X}^2-\mathbb{\tilde{X}}^2\vert_{C_2^{2\nu}}+\vert Y'-\tilde{Y}'\vert_{C^{\nu}}+\vert R^Y-R^{\tilde{Y}}\vert_{C_2^{2\nu}}+\vert Y-\tilde{Y}\vert_{C^{\nu}}),
\end{multline}
where
$$
C_X=1+\vert X\vert_{C^{\nu}}+\vert \tilde{X}\vert_{C^{\nu}}+\vert \mathbb{X}^2\vert_{C_2^{2\nu}}+\vert \mathbb{\tilde{X}}^2\vert_{C_2^{2\nu}}.
$$

In the case of $X=\tilde{X}$, inequality \eqref{eqn:ParametricRoughPath-2general} can be rewritten as
\begin{equation}
\Vert \nabla^nV^Y-\nabla^nV^{\tilde{Y}}\Vert_{L^{\infty}}\leq
16C_1C_2^3\Vert \nabla^{n+1}A\Vert_{C^2}\Vert Y\Vert_{\mathcal{D}_X}(1+\Vert Y\Vert_{\mathcal{D}_X})^2\Vert Y-\tilde{Y}\Vert_{\mathcal{D}_X}^*.
\label{eqn:ParametricRoughPath-2}
\end{equation}
\end{lemma}
\begin{proof}[ Proof of Lemma \ref{lem:ParametricRoughPath}]
Inequalities \eqref{eqn:ParametricRoughPath-1} and
\eqref{eqn:ParametricRoughPath-2} were proved in
\cite{[Bess+Gub_2005]}, Lemma $7$. Now we will show
\eqref{eqn:ParametricRoughPath-2general}. It is enough to consider
the case of $n=0$. By formulae \eqref{eqn:IntegralDef-3} and
\eqref{eqn:IntegralDef-4} we have
\begin{eqnarray}
V^Y-V^{\tilde{Y}}(x)&=&A(x-Y(0))(Y(1)-Y(0))
\nonumber\\
&-&A(x-\tilde{Y}(0))(\tilde{Y}(1)-\tilde{Y}(0))\nonumber\\
&+&(A(x-Y))'(0)Y'(0)\mathbb{X}^2(0,1)
\nonumber\\
&-&(A(x-\tilde{Y}))'(0)\tilde{Y}'(0)\mathbb{\tilde{X}}^2(0,1)
\nonumber\\
&+&Q^x(0,1)-\tilde{Q}^x(0,1)\nonumber
\end{eqnarray}
where $Q^x$ and $\tilde{Q}^x$ (given by formulae
\eqref{eqn:IntegralDef-3} and \eqref{eqn:IntegralDef-4}) satisfy
inequality \eqref{eqn:IntegralDef-5} and we have identified $S^1$
with $[0,1]$. Therefore, $Y(1)=Y(0)$, $\tilde{Y}(1)=\tilde{Y}(0)$.
Hence, we have
\begin{multline}\label{eqn:PRaux-1}
\vert V^Y-V^{\tilde{Y}}\vert_{L^{\infty}}\leq
\supl_x\vert (A(x-Y))'(0)Y'(0)\mathbb{X}^2(0,1)-
(A(x-\tilde{Y}))'(0)\tilde{Y}'(0)\mathbb{\tilde{X}}^2(0,1)\vert \\
+\supl_x\vert Q^x(0,1)-\tilde{Q}^x(0,1)\vert .
\end{multline}
For the first term on the r.h.s. we have
\begin{align}\label{eqn:PRaux-2}
\vert (&A(x-Y))'(0)Y'(0)\mathbb{X}^2(0,1)-
(A(x-\tilde{Y}))'(0)\tilde{Y}'(0)\mathbb{\tilde{X}}^2(0,1)\vert \\
&\leq
\vert (dA(x-Y(0))Y'(0)Y'(0)-dA(x-\tilde{Y}(0))\tilde{Y}'(0)\tilde{Y}'(0))\mathbb{X}^2(0,1)\vert \nonumber\\
&+\vert dA(x-\tilde{Y}(0))\tilde{Y}'(0)\tilde{Y}'(0)(\mathbb{X}^2(0,1)-\mathbb{\tilde{X}}^2(0,1))\vert \nonumber\\
&\leq \vert \mathbb{X}^2\vert_{C_2^{2\nu}}\vert dA(x-Y(0))Y'(0)Y'(0)-dA(x-\tilde{Y}(0))\tilde{Y}'(0)\tilde{Y}'(0)\vert \nonumber\\
&+ \vert A\vert_{C^2}\vert Y'\vert_{L^{\infty}}^2\vert \mathbb{X}^2-\mathbb{\tilde{X}}^2\vert_{C_2^{2\nu}}\nonumber\\
&\leq
\vert \mathbb{X}^2\vert_{C_2^{2\nu}}\vert Y'\vert_{L^{\infty}}^2\vert A\vert_{C^2}\vert Y'-\tilde{Y}'\vert_{L^{\infty}}
+\vert \mathbb{X}^2\vert_{C_2^{2\nu}}\vert A\vert_{C^1}(\vert Y'\vert_{L^{\infty}}+\vert \tilde{Y}'\vert_{L^{\infty}})\vert Y'-\tilde{Y}'\vert_{L^{\infty}}\nonumber\\
&+\vert A\vert_{C^2}\vert Y'\vert_{L^{\infty}}^2\vert \mathbb{X}^2-\mathbb{\tilde{X}}^2\vert_{C_2^{2\nu}}.\quad\nonumber
\end{align}
By \eqref{eqn:IntegralDef-5} we can estimate second term as
follows
\begin{align}
\vert Q^x-\tilde{Q}^x\vert_{C_2^{3\nu}}&\leq
C\Big[(\Vert A(x-Y)\Vert_{\mathcal{D}_X}+\Vert A(x-\tilde{Y})\Vert_{\mathcal{D}_{\tilde{X}}})\eps_Y\nonumber\\
&+
(\Vert Y\Vert_{\mathcal{D}_X}+\Vert \tilde{Y}\Vert_{\mathcal{D}_{\tilde{X}}})\eps_A+\eps_X\Big].\label{eqn:PRaux-3}
\end{align}
where
\begin{align*}
\eps_Y&=\vert Y-\tilde{Y}\vert_{C^{\nu}}+\vert Y'-\tilde{Y}'\vert_{C^{\nu}}+\vert R^Y-R^{\tilde{Y}}\vert_{C_2^{2\nu}},\\
\eps_A&=\vert A(x-Y)-A(x-\tilde{Y})\vert_{C^{\nu}}+\vert A(x-Y)'-A(x-\tilde{Y})'\vert_{C^{\nu}}\\
&+\vert R^{A(x-Y)}-R^{A(x-\tilde{Y})}\vert_{C_2^{2\nu}},\\
\eps_X&=(\Vert A(x-Y)\Vert_{\mathcal{D}_X}+\Vert A(x-\tilde{Y})\Vert_{\mathcal{D}_{\tilde{X}}})\\
&\times(\Vert Y\Vert_{\mathcal{D}_X}+\Vert \tilde{Y}\Vert_{\mathcal{D}_{\tilde{X}}})
(\vert X-\tilde{X}\vert_{C^{\nu}}+\vert \mathbb{X}^2-\mathbb{\tilde{X}}^2\vert_{C_2^{2\nu}}).
\end{align*}
By formula \eqref{eqn:RegularMap-4} we can estimate $\eps_A$ as
follows
\begin{align}
\vert \eps_A\vert &\leq
K\vert A\vert_{C^3}(1+\vert X\vert_{C^{\nu}}+\vert \tilde{X}\vert_{C^{\nu}})^3(1+\vert Y\vert_{D_X}+\vert \tilde{Y}\vert_{D_{\tilde{X}}})^2\nonumber\\
&\times(\vert X-\tilde{X}\vert_{C^{\nu}}+\vert Y-\tilde{Y}\vert_{C^{\nu}}+\vert Y'-\tilde{Y}'\vert_{C^{\nu}}+\vert R^Y-R^{\tilde{Y}}\vert_{C_2^{2\nu}}).\label{eqn:PRaux-4}
\end{align}
By inequality \eqref{eqn:RegularMap-3} we infer that
\begin{eqnarray}
\Vert A(x-Y)\Vert_{\mathcal{D}_X}\leq
K\vert A\vert_{C^2}\vert Y\vert_{D_X}(1+\vert Y\vert_{D_X})(1+\vert X\vert_{C^{\nu}})^2,\label{eqn:PRaux-5}
\end{eqnarray}
and similarly,
\begin{eqnarray}
\Vert A(x-\tilde{Y})\Vert_{\mathcal{D}_{\tilde{X}}}\leq
K\vert A\vert_{C^2}\vert \tilde{Y}\vert_{D_{\tilde{X}}}(1+\vert \tilde{Y}\vert_{D_{\tilde{X}}})(1+\vert \tilde{X}\vert_{C^{\nu}})^2.\label{eqn:PRaux-6}
\end{eqnarray}
Therefore, combining \eqref{eqn:PRaux-3} with \eqref{eqn:PRaux-4},
\eqref{eqn:PRaux-5} and \eqref{eqn:PRaux-6} we get
\begin{multline}\label{eqn:PRaux-7}
\vert Q^x-\tilde{Q}^x\vert_{C_2^{3\nu}}\leq C(\nu)\vert A\vert_{C^{n+3}}(1+\vert X\vert_{C^{\nu}}+\vert \tilde{X}\vert_{C^{\nu}})^4(1+\Vert Y\Vert_{\mathcal{D}_X}+\Vert \tilde{Y}\Vert_{\mathcal{D}_{\tilde{X}}})^3\\
(\vert X-\tilde{X}\vert_{C^{\nu}}+\vert \mathbb{X}^2-\mathbb{\tilde{X}}^2\vert_{C_2^{2\nu}}+\vert Y'-\tilde{Y}'\vert_{C^{\nu}}+\vert R^Y-R^{\tilde{Y}}\vert_{C_2^{2\nu}}+\vert Y-\tilde{Y}\vert_{C^{\nu}}).
\end{multline}
Hence, the result follows from \eqref{eqn:PRaux-2} and
\eqref{eqn:PRaux-7}.
\end{proof}
We will denote for any $Y\in\mathcal{D}_X$,
$\tilde{Y}\in\mathcal{D}_{\tilde{X}}$
$$
\vert Y-\tilde{Y}\vert_D=\vert X-\tilde{X}\vert_{C^{\nu}}+\vert \mathbb{X}^2-\mathbb{\tilde{X}}^2\vert_{C_2^{2\nu}}+\vert Y'-\tilde{Y}'\vert_{C^{\nu}}+\vert R^Y-R^{\tilde{Y}}\vert_{C_2^{2\nu}}+\vert Y-\tilde{Y}\vert_{C^{\nu}}.
$$

\section{Random Filaments evolution problem}
Let $\mathcal{D}_{\mathbb{X},T}=C([0,T],\mathcal{D}_X)$ be a
vector space with the usual supremum norm
\begin{equation}
\vert F\vert_{\mathcal{D}_{\mathbb{X},T}}=\supl_{t\in[0,T]}\vert F(t)\vert_{\mathcal{D}_X}^*\label{eqn:NormDef}.
\end{equation}
Obviously $\mathcal{D}_{\mathbb{X},T}$ is a Banach space. Assume
also that the function $\phi$ appeared in the formula
\eqref{eqn:RandomFilam-3} satisfies following hypothesis.
\begin{hypothesis}\label{hyp:PhiAssumptions}
\begin{trivlist}
\item[(i)]$\phi:\mathbb{R}^3\to\mathbb{R}$ is even function.
\item[(ii)] the Fourier transform of $\phi$ is real and
non-negative function:
$$
\hat{\phi}(k)\geq 0,\;\;k\in\Rnu^3
$$
\item[(iii)]
$$
\intl_{\Rnu^3}(1+\vert k\vert^2)^2\hat{\phi}(k)dk<\infty
$$
\end{trivlist}
\end{hypothesis}
\begin{example}
The function $\phi=\phi_\mu$, $\mu>0$ defined by
$$\phi(\cdot)=\frac{\Gamma}{(\vert \cdot\vert^2+\mu^2)^{1/2}}$$
is smooth and satisfies hypothesis \ref{hyp:PhiAssumptions}, see p.6
of \cite{[Bers+Gub_2007]}.
\end{example}
Then the following local existence and uniqueness Theorem for
problem \eqref{eqn:RandomFilam-1}-\eqref{eqn:RandomFilam-3} has
been proved in \cite{[Bess+Gub_2005]}, see Theorem
$3$,p.1842.
\begin{theorem}\label{thm:LocalExistUniq}
Assume $\phi\in C^6(\Rnu^3,\Rnu)$, $\nu\in (\frac{1}{3},1)$,
$\mathbb{X}=(X,\mathbb{X}^2)$ is a $\nu$-rough path, $\gamma_0\in
\mathcal{D}_X$. Then there exists a time
$T_0=T_0(\nu,\vert \phi\vert_{C^5},\mathbb{X})>0$ such that the problem
\eqref{eqn:RandomFilam-1}-\eqref{eqn:RandomFilam-3} has unique
solution in $\mathcal{D}_{\mathbb{\gamma_0},T_0}\subset\mathcal{D}_{\mathbb{X},T_0}$.
\end{theorem}
Our aim is to prove global existence of solution of the problem
\eqref{eqn:RandomFilam-1}-\eqref{eqn:RandomFilam-3} under
assumptions of Theorem \ref{thm:LocalExistUniq} and additional
hypothesis \ref{hyp:PhiAssumptions} i.e. we shall prove
\begin{theorem}\label{thm:GlobalExistUniq}
Assume that $T>0$, $\gamma_0$ is a geometric $\nu$-rough path, $\phi\in
C^6(\Rnu^3,\Rnu)$ satisfies  satisfies Hypothesis
\ref{hyp:PhiAssumptions} and  $\nu\in(\frac{1}{3},1)$. Then the problem
\eqref{eqn:RandomFilam-1}-\eqref{eqn:RandomFilam-3} has unique
solution in $\mathcal{D}_{\mathbb{\gamma}_0,T}$.
\end{theorem}
We will need the following definition.
\begin{definition}
Let $\phi\in C^4(\Rnu^3,\Rnu)$, $\gamma\in \mathcal{D}_X$. Put
\begin{equation}
\mathcal{H}_X^{\phi}(\gamma)=\frac{1}{2}\intl_{S^1}\vec{\psi}^{\gamma}(\gamma(\xi))\cdot
d\vec{\gamma}(\xi),\label{eqn:EnergyIntegralDef-1}
\end{equation}
where
\begin{equation*}
\vec{\psi}^{\gamma}(x)=\intl_{S^1}\phi(x-\gamma(\eta))d\vec{\gamma}(\eta).\label{eqn:EnergyIntegralDef-1a}
\end{equation*}
$\mathcal{H}_X^{\phi}(\gamma)$ is called the energy of path $\gamma$. We will omit $\phi$ below.
\end{definition}
\begin{remark}\label{rem:CorrectEnergyDef}
Definition
\eqref{eqn:EnergyIntegralDef-1}-\eqref{eqn:EnergyIntegralDef-1a}
is well posed. Indeed, by Lemma \ref{lem:ParametricRoughPath}
$\vec{\psi}^{\gamma}\in C^2(\Rnu^3,\Rnu^3)$ and, therefore, it
follows by Lemma \ref{lem:RegularMap} that
$\vec{\psi}^{\gamma}(\gamma(\cdot))\in \mathcal{D}_X$.
\end{remark}
\begin{remark}
Assume that $\nu>\frac{1}{2}$ and $\gamma\in C^1(S^1,\Rnu^3)$.
Then by Remark \ref{rem:SmoothCase} the line integrals in the
definition of energy are understood in the sense of Young and
\begin{equation}
\mathcal{H}_X(\gamma)=\frac{1}{2}\intl_{S^1}\intl_{S^1}\phi(\vec{\gamma}(\xi)-\vec{\gamma}(\eta))
(\frac{d\vec{\gamma}}{d\xi}(\xi),\frac{d\vec{\gamma}}{d\eta}(\eta))d\xi
d\eta. \label{eqn:EnergyIntegralDef-2}
\end{equation}
\end{remark}
\begin{lemma}\label{lem:EnergyProperties}
Assume $\phi\in C^4(\Rnu^3,\Rnu)$. Then there exists constant
$C=C(\nu,\mathbb{X})$ such that for all $\gamma\in \mathcal{D}_X$
\begin{equation}
\vert \mathcal{H}_X(\gamma)\vert \leq
C\vert \phi\vert_{C^4}\vert \gamma\vert_{\mathcal{D}_X}^4(1+\vert \gamma\vert_{\mathcal{D}_X})^2.\label{eqn:EnergyEstimAbove}
\end{equation}
Moreover, the map $\mathcal{H}_X:\mathcal{D}_X\to\Rnu$ is locally
Lipshitz i.e. for any $R>0$ there exists $C=C(R)$ such that for
any $\gamma,\tilde{\gamma}\in\mathcal{D}_X$,
$\vert \gamma\vert_{\mathcal{D}_X}<R$, $\vert \tilde{\gamma}\vert_{\mathcal{D}_X}<R$
we have
\begin{equation}
\vert \mathcal{H}_X(\gamma)-\mathcal{H}_X(\tilde{\gamma})\vert \leq
C(R)\vert \gamma-\tilde{\gamma}\vert_{\mathcal{D}_X}^*.\label{eqn:LocLipshitz}
\end{equation}
Furthermore, for any $R>0$ there exists $C=C(R)$ such that for any
$\gamma\in\mathcal{D}_X$,
$\tilde{\gamma}\in\mathcal{D}_{\tilde{X}}$,
$$
\vert \gamma\vert_{\mathcal{D}_X}<R,
\vert \tilde{\gamma}\vert_{\mathcal{D}_{\tilde{X}}}<R,
$$
$$
C_{X,\tilde{X}}=\vert X\vert_{C^{\nu}}+\vert \tilde{X}\vert_{C^{\nu}}+\vert \mathbb{X}^2\vert_{C_2^{2\nu}}+\vert \mathbb{\tilde{X}}^2\vert_{C_2^{2\nu}}<R
$$
we have
\begin{equation}
\vert \mathcal{H}_X(\gamma)-\mathcal{H}_{\tilde{X}}(\tilde{\gamma})\vert \leq
C(R)(\vert X-\tilde{X}\vert_{C^{\nu}}+\vert \mathbb{X}^2-\mathbb{\tilde{X}}^2\vert_{C_2^{2\nu}}+\vert \gamma-\tilde{\gamma}\vert_{C^{\nu}}+\vert \gamma'-\tilde{\gamma}'\vert_{C^{\nu}}+\vert R^{\gamma}-R^{\tilde{\gamma}}\vert_{C_2^{2\nu}}).\label{eqn:LocLipshitz-2}
\end{equation}
\end{lemma}
\begin{proof}[ Proof of Lemma \ref{lem:EnergyProperties}]
First we will show inequality \eqref{eqn:EnergyEstimAbove}. By
representation \eqref{eqn:IntegralDef-2} we have
\begin{align}
\mathcal{H}_X(\gamma)&=\frac{1}{2}(\vec{\psi}^{\gamma}(\gamma(0))(\vec{\gamma}(1)-\vec{\gamma}(0)))\nonumber\\
&+\left[d\vec{\psi}^{\gamma}(\gamma(0))\gamma'(0)\right]\gamma'(0)\mathbb{X}^2(1,0)+Q(0,1)\nonumber\\
&=I+II+III,\gamma\in\mathcal{D}_X.\nonumber
\end{align}
Since $\vec{\gamma}(1)=\vec{\gamma}(0)$ we infer that $I=0$.
Concerning the second term by Lemma \ref{lem:ParametricRoughPath}
we have the following estimate
\begin{align}
\vert II\vert &\leq
\Vert \mathbb{X}^2\Vert_{C_2^{2\nu}}\Vert \nabla\vec{\psi}\Vert_{L^{\infty}}\vert \gamma'\vert_{L^{\infty}}^2\nonumber\\
&\leq
C(\nu,\mathbb{X})\vert \phi\vert_{C^3}\vert \gamma\vert_{\mathcal{D}_X}^4(1+\vert \gamma\vert_{\mathcal{D}_X}).\label{eqn:EnergyEstimAbove-1}
\end{align}
For third term we infer from inequality \eqref{eqn:IntegralDef-3}
\begin{equation}
\vert III\vert \leq \Vert Q\Vert_{C_2^{3\nu}}\leq
C(\nu,\mathbb{X})\vert \vec{\psi}(\vec{\gamma})\vert_{\mathcal{D}_X}\vert \vec{\gamma}\vert_{\mathcal{D}_X}.\label{eqn:EnergyEstimAbove-2}
\end{equation}
Then by Lemmas \ref{lem:RegularMap} and
\ref{lem:ParametricRoughPath} we have
\begin{align}
\vert \vec{\psi}^{\gamma}(\vec{\gamma})\vert_{\mathcal{D}_X}&\leq
C(\nu,\mathbb{X})\vert \psi^{\gamma}\vert_{C^2}\vert \vec{\gamma}\vert_{\mathcal{D}_X}(1+\vert \vec{\gamma}\vert_{\mathcal{D}_X})\nonumber\\
&\leq C(\nu,\mathbb{X})\vert \phi\vert_{C^4}\vert \vec{\gamma}\vert_{\mathcal{D}_X}^3(1+\vert \vec{\gamma}\vert_{\mathcal{D}_X})^2\label{eqn:EnergyEstimAbove-3}
\end{align}
Combining \eqref{eqn:EnergyEstimAbove-1},
\eqref{eqn:EnergyEstimAbove-2} and \eqref{eqn:EnergyEstimAbove-3}
we get inequality \eqref{eqn:EnergyEstimAbove}.

Now we will prove inequality \eqref{eqn:LocLipshitz-2}. By formula
\eqref{eqn:IntegralDef-4} we have
\begin{align}\label{eqn:EnergyLocLip-1}
\mathcal{H}_X(\gamma)-\mathcal{H}_{\tilde{X}}(\tilde{\gamma})&=\frac{1}{2}
\Big[(\nabla\psi^{\gamma}(\gamma(0))\gamma'(0)\gamma'(0)-\nabla\psi^{\tilde{\gamma}}(\tilde{\gamma}(0))\tilde{\gamma}'(0)\tilde{\gamma}'(0))
\mathbb{X}^2(1,0)\nonumber\\
&+\nabla\psi^{\tilde{\gamma}}(\tilde{\gamma}(0))\tilde{\gamma}'(0)\tilde{\gamma}'(0)(\mathbb{X}^2(1,0)-\mathbb{\tilde{X}}^2(1,0))
+Q(0,1)-\tilde{Q}(0,1)\Big]\!
\nonumber\\
 &=\!I\! +\! II\! +\! III
\end{align}
The first term in \eqref{eqn:EnergyLocLip-1} can be represented as
follows
\begin{eqnarray}
I &=&
(\nabla\psi^{\gamma}(\gamma(0))\gamma'(0)\gamma'(0)-\nabla\psi^{\tilde{\gamma}}(\tilde{\gamma}(0))\tilde{\gamma}'(0)\tilde{\gamma}'(0))
\mathbb{X}^2(1,0)\nonumber\\
 &=& \left[(\nabla\psi^{\gamma}(\gamma(0))-\nabla\psi^{\tilde{\gamma}}(\tilde{\gamma}(0)))\gamma'(0)\gamma'(0)\right.\nonumber\\
 &+& \nabla\psi^{\tilde{\gamma}}(\tilde{\gamma}(0))(\gamma'(0)-\tilde{\gamma}(0))\gamma'(0)\nonumber\\
 &+& \left.\nabla\psi^{\tilde{\gamma}}(\tilde{\gamma}(0))\tilde{\gamma}(0)(\gamma'(0)-\tilde{\gamma}(0))\right]\mathbb{X}^2(1,0)=A+B+C\label{eqn:EnergyLocLip-2}
\end{eqnarray}
The first term in \eqref{eqn:EnergyLocLip-2} can be estimated as
follows
\begin{eqnarray}\label{eqn:EnergyLocLip-3}
\vert A\vert &=&\vert (\nabla\psi^{\gamma}(\gamma(0))-\nabla\psi^{\tilde{\gamma}}(\tilde{\gamma}(0)))\gamma'(0)\gamma'(0)\mathbb{X}^2(1,0)\vert \\
&\leq&
\Vert \mathbb{X}^2\Vert_{C_2^{2\nu}}\vert \gamma\vert_{\mathcal{D}_X}^2(\vert \nabla\psi^{\gamma}
(\gamma(0))\nonumber\\
&-&\nabla\psi^{\gamma}(\tilde{\gamma}(0))\vert +\vert \nabla\psi^{\gamma}(\tilde{\gamma}(0))-\nabla\psi^{\tilde{\gamma}}(\tilde{\gamma}(0))\vert )\nonumber\\
&\leq &\Vert \mathbb{X}^2\Vert_{C_2^{2\nu}}\vert \gamma\vert_{\mathcal{D}_X}^2(\vert \psi^{\gamma}\vert_{C^2}\vert \gamma(0)
\nonumber\\
&-&\tilde{\gamma}(0)\vert
+C_X^4\vert \phi\vert_{C^4}(1+\vert \gamma\vert_{\mathcal{D}_X}+\vert \tilde{\gamma}\vert_{\mathcal{D}_{\tilde{X}}})^3\vert \gamma-\tilde{\gamma}\vert_D
\nonumber\\
&\leq&
KC_X^4\vert \phi\vert_{C^4}(1+\vert \gamma\vert_{\mathcal{D}_X}+\vert \tilde{\gamma}\vert_{\mathcal{D}_{\tilde{X}}})^3\vert \gamma-\tilde{\gamma}\vert_{D},
\nonumber
\end{eqnarray}
where second inequality follows from inequality
\eqref{eqn:ParametricRoughPath-2general} and third one from
inequality \eqref{eqn:ParametricRoughPath-1}. For second term we
have by inequality \eqref{eqn:ParametricRoughPath-1}
\begin{eqnarray}
\vert B\vert &\leq&
C\Vert \mathbb{X}^2\Vert_{C_2^{2\nu}}\vert \gamma\vert_{\mathcal{D}_X}\vert \phi\vert_{C^3}\vert \tilde{\gamma}\vert_{\mathcal{D}_X}^2
(1+\vert \tilde{\gamma}\vert_{\mathcal{D}_X})\vert \gamma-\tilde{\gamma}\vert_{D}
\nonumber
\\
&\leq& CC_X(1+\vert \gamma\vert_{\mathcal{D}_X}+\vert \tilde{\gamma}\vert_{\mathcal{D}_{\tilde{X}}})^3\vert \gamma-\tilde{\gamma}\vert_D.
\label{eqn:EnergyLocLip-4}
\end{eqnarray}
Similarly, we have for third term
\begin{eqnarray}
\vert C\vert \leq
C(\nu,\mathbb{X},\vert \gamma\vert_{\mathcal{D}_X},\vert \tilde{\gamma}\vert_{\mathcal{D}_X})\vert \gamma-\tilde{\gamma}\vert_{D}.\label{eqn:EnergyLocLip-5}
\end{eqnarray}
Term $II$ in \eqref{eqn:EnergyLocLip-1} can be estimated as
follows
\begin{align*}
\vert II\vert &\leq
\vert \nabla\psi^{\tilde{\gamma}}\vert_{L^{\infty}}\vert \tilde{\gamma}\vert_{\mathcal{D}_X}^2\vert \gamma-\tilde{\gamma}\vert_D\\
&\leq
C_X^3\vert \phi\vert_{C^3}(1+\vert \gamma\vert_{\mathcal{D}_X}+\vert \tilde{\gamma}\vert_{\mathcal{D}_{\tilde{X}}})^3\vert \gamma-\tilde{\gamma}\vert_D.
\end{align*}
Thus it remains to estimate third term of equality
\eqref{eqn:EnergyLocLip-1}. We have by inequality
\eqref{eqn:IntegralDef-5}
\begin{align*}
\vert Q(0,1)-\tilde{Q}(0,1)\vert &\leq \Vert Q-\tilde{Q}\Vert_{C_2^{3\nu}}\\
&\leq
C_X\Big[(\vert \psi^{\tilde{\gamma}}(\tilde{\gamma})\vert_{\mathcal{D}_{\tilde{X}}}+
\vert \psi^{\gamma}(\gamma)\vert_{\mathcal{D}_X})\vert \gamma-\tilde{\gamma}\vert_D\\
&+(\vert \tilde{\gamma}\vert_{\mathcal{D}_{\tilde{X}}}+\vert \tilde{\gamma}\vert_{\mathcal{D}_{\tilde{X}}})\vert \psi^{\tilde{\gamma}}(\tilde{\gamma})-\psi^{\gamma}(\gamma)\vert_D\\
&+(\vert \psi^{\tilde{\gamma}}(\tilde{\gamma})\vert_{\mathcal{D}_{\tilde{X}}}+
\vert \psi^{\gamma}(\gamma)\vert_{\mathcal{D}_X})(\vert \tilde{\gamma}\vert_{\mathcal{D}_{\tilde{X}}}+\vert \tilde{\gamma}\vert_{\mathcal{D}_{\tilde{X}}})\\
&\times(\vert X-\tilde{X}\vert_{C^{\nu}}+\vert \mathbb{X}^2-\mathbb{\tilde{X}}^2\vert_{C_2^{2\nu}})\Big]
\end{align*}
Term $\vert \psi^{\tilde{\gamma}}(\tilde{\gamma})\vert_{\mathcal{D}_X}$ is
bounded by the constant
$C=C(\nu,\mathbb{X},\vert \tilde{\gamma}\vert_{\mathcal{D}_X})$ by
inequality \eqref{eqn:ParametricRoughPath-1}. Therefore, to prove
estimate \eqref{eqn:LocLipshitz} it is enough to show that there
exists constant $C=C(\nu,\mathbb{X},R)$ such that for
$\gamma,\tilde{\gamma}\in\mathcal{D}_X$ with
$\vert \gamma\vert_{\mathcal{D}_X},\vert \tilde{\gamma}\vert_{\mathcal{D}_X}\leq R$
\begin{equation}
\vert \psi^{\tilde{\gamma}}(\tilde{\gamma})-\psi^{\gamma}(\gamma)\vert_D\leq
C\vert \tilde{\gamma}-\gamma\vert_D.\label{eqn:EnergyLocLip-6}
\end{equation}
By triangle inequality we have
\begin{align}
\vert \psi^{\tilde{\gamma}}(\tilde{\gamma})-\psi^{\gamma}(\gamma)\vert_D
&\leq\vert \psi^{\tilde{\gamma}}(\tilde{\gamma})-\psi^{\tilde{\gamma}}(\gamma)\vert_D+
\vert \psi^{\tilde{\gamma}}(\gamma)-\psi^{\gamma}(\gamma)\vert_D\nonumber\\
&=I+II.\label{eqn:EnergyLocLip-7}
\end{align}
The first term can be estimated using inequality
\eqref{eqn:RegularMap-4} as follows
\begin{equation}
\vert I\vert \leq
KC_X^3\vert \psi^{\tilde{\gamma}}\vert_{C^3}(1+\vert \tilde{\gamma}\vert_{\mathcal{D}_{\tilde{X}}}+
\vert \gamma\vert_{\mathcal{D}_X})^2\vert \tilde{\gamma}-\gamma\vert_D.\label{eqn:EnergyLocLip-8}
\end{equation}
By inequality \eqref{eqn:ParametricRoughPath-1} we have
\begin{equation}
\vert \psi^{\tilde{\gamma}}\vert_{C^3}\leq
C\vert \phi\vert_{C^5}\vert \tilde{\gamma}\vert_{\mathcal{D}_X}^2(1+\vert \tilde{\gamma}\vert_{\mathcal{D}_X})\label{eqn:EnergyLocLip-8'}
\end{equation}
Combining \eqref{eqn:EnergyLocLip-8} and
\eqref{eqn:EnergyLocLip-8'} we get necessary estimate for $I$. It
remains to find an estimate for term $II$. By inequalities
\eqref{eqn:RegularMap-3} and \eqref{eqn:ParametricRoughPath-2} we
have
\begin{align}
II&=\vert \psi^{\tilde{\gamma}}(\gamma)-\psi^{\gamma}(\gamma)\vert_D\nonumber\\
&\leq
(1+\vert X\vert_{\nu})\vert \psi^{\tilde{\gamma}}(\gamma)-\psi^{\gamma}(\gamma)\vert_{\mathcal{D}_X}\nonumber\\
&\leq K\vert \nabla\psi^{\tilde{\gamma}}-\nabla\psi^{\gamma}\vert_{C^1}
\vert \gamma\vert_{\mathcal{D}_X}(1+\vert \gamma\vert_{\mathcal{D}_X})(1+\vert X\vert_{\nu})^3
\nonumber\\
&\leq K\vert \phi\vert_{C^5}C_X^7(1+\vert \gamma\vert_{\mathcal{D}_X}+\vert \tilde{\gamma}\vert_{\mathcal{D}_{\tilde{X}}})^5
\vert \tilde{\gamma}-\gamma\vert_D.\label{eqn:EnergyLocLip-9}
\end{align}
Hence the inequality \eqref{eqn:LocLipshitz-2} follows. Inequality
\eqref{eqn:LocLipshitz} is a consequence of inequality
\eqref{eqn:LocLipshitz-2}.
\end{proof}
\begin{corollary}\label{cor:ContinuityEnergy}
Under assumptions of Lemma \ref{lem:EnergyProperties} and
assumption \ref{ass:RoughPathApproximability} energy
$\mathcal{H}_X:\mathcal{D}_X\to\Rnu$ is a continuous function on
$\mathcal{D}_X$. Furthermore, for any $\gamma\in\mathcal{D}_X$
$$
0\leq\mathcal{H}_X(\gamma)<\infty.
$$
\end{corollary}
\begin{proof}[Proof of Corollary \ref{cor:ContinuityEnergy}]
We only need to show that $\mathcal{H}_X(\gamma)\geq 0$, for any
$\gamma\in\mathcal{D}_X$. Other statements of the Corollary easily
follow from Lemma \ref{lem:EnergyProperties}.
%It is enough to show the following formula
Fix $n\in \mathbb{N}$.
Let $C(0)=\cup_{i=1}^{n^6}C(k_i)^n$ be a partition of the cube $C(0)$ with center $0$ and length $n^2$ on the cubes $C(k_i)^n$ of the length of $\frac{1}{n}$ with centers $k_i$ and nonintersecting interiors. Define $\hat{\phi}^n(k)=\hat{\phi}(k_i)$, $k\in C(k_i)^n$ and $0$ otherwise. Consequently, define
\[
\phi^n(x)=\suml_{i=1}^{n^6}\hat{\phi}(k_i)\int\limits_{C(k_i)^n}e^{i<k,x>}dk,x\in\mathbb{R}^3
\]
Then $\phi=\lim\limits_{n\to\infty}\phi^n$ in $C^4$ norm because of the assumption \ref{hyp:PhiAssumptions}.
Consequently,
\begin{equation}
\mathcal{H}_X^{\phi}(\gamma)=\lim\limits_{n\to\infty}\mathcal{H}_X^{\phi_n}(\gamma).\label{aux:1}
\end{equation}
Now formula
\begin{equation}
\mathcal{H}_X^{z}(\gamma)=\frac{1}{(2\pi)^3}\intl_{\Rnu^3}\hat{z}(k)\Big\vert \intl_{S^1}e^{i(k,\gamma_s)}d\gamma_s\Big\vert^2dk,\gamma\in
\mathcal{D}_X.\label{eqn:EnergyIntegralRepr-1}
\end{equation}
is correct for $z=\phi_n$, because $\phi^n$ is a sum of a finite number of Fourier modes. Therefore,
\begin{equation}
\mathcal{H}_X^{\phi_n}(\gamma)\geq 0.\label{aux:2}
\end{equation}
Thus, the  result follows from identity \eqref{aux:1} and inequality \eqref{aux:2}.
\end{proof}
Now we are going to show that energy is a local integral of motion
for problem \eqref{eqn:RandomFilam-1}-\eqref{eqn:RandomFilam-3}.
\begin{lemma}\label{lem:EnergyMotionIntegral}
Let $\gamma\in\mathcal{D}_{\mathbb{\gamma}_0,T_0}$ be a local solution of
problem \eqref{eqn:RandomFilam-1}-\eqref{eqn:RandomFilam-3} (such
solution exists by Theorem \ref{thm:LocalExistUniq}). Then
$$
\frac{d\mathcal{H}_{\gamma_0}(\gamma(s))}{ds}=0,s\in[0,T_0).
$$
\end{lemma}
\begin{proof}[Proof of Lemma \ref{lem:EnergyMotionIntegral}]
Since $\gamma(0)=\gamma_0\in \mathcal{D}_{\gamma_0}$ is a geometric rough path
there exist sequence $\{\gamma_0^n\}_{n=1}^{\infty}\in C^{\infty}(S^1,\Rnu^3)$ such that
\[
\vert \gamma_0^n-\gamma_0\vert_{C^{\nu}}+\vert \mathbb{\gamma}^2-\mathbb{\gamma}^2\vert_{C_2^{2\nu}}\to 0,n\to\infty.
\]
Now $\gamma_0^n\in \mathcal{D}_{\gamma_0^n}$, $\gamma_0\in \mathcal{D}_{\gamma_0}$. Thus we can put
$(\gamma_0^n)'=(\gamma_0)'=1$ and $R^{\gamma_0^n}=R^{\gamma_0}=0$.
Hence we deduce that
$$
\vert \gamma_0^n-\gamma_0\vert_D\to 0,n\to\infty.
$$
Denote $\gamma^n\in C([0,\infty),\mathbf{H}^1(S^1,\Rnu^3))$ the
global solution of problem
\eqref{eqn:RandomFilam-1}-\eqref{eqn:RandomFilam-3} with initial
condition $\gamma_0^n$. Existence of such solution has been proved in
Theorem $2$ of \cite{[Bers+Gub_2007]}. Then according to
\cite{[Bers+Gub_2007]} (Theorem $4$, p.1846) we have
$$
\supl_{t\in[0,T_0]}\vert \gamma^n(t)-\gamma(t)\vert_D\to 0,n\to\infty.
$$
Therefore, by continuity of energy functional $\mathcal{H}_{\gamma_0}$ we
have
\begin{equation}
\mathcal{H}_{\gamma_0}(\gamma(s))=\liml_{n\to\infty}\mathcal{H}_{\gamma_0^n}(\gamma^n(s)),s\in[0,T_0].\label{eqn:lemEMIaux-1}
\end{equation}
Furthermore, by Lemma $2$ of \cite{[Bers+Gub_2007]}, we
have
\begin{equation}
\mathcal{H}_{\gamma_0^n}(\gamma^n(s))=\mathcal{H}_{\gamma_0^n}(\gamma_0^n),s\in[0,T_0].\label{eqn:lemEMIaux-2}
\end{equation}
As a result, combining \eqref{eqn:lemEMIaux-1} and
\eqref{eqn:lemEMIaux-2} we get statement of the Lemma.
\end{proof}
Let us recall definition of along with $\nu$-rough path $Y$
\begin{equation}
u^Y(x)=\intl_Y\nabla\phi(x-y)\times dy,Y\in
\mathcal{D}_X.\label{eqn:RandomFilam-3'}
\end{equation}
Now we show that if energy functional of $Y$ is bounded then
associated velocity is a smooth function. We have
\begin{lemma}[See Lemma $3$ in \cite{[Bers+Gub_2007]}]\label{lem:VelocitySmoothness}
For any $n\in \mathbb{Z},n\geq 0$, we have following bound
\begin{equation}
\Vert \nabla^n u^{\gamma}\Vert_{L^{\infty}}\leq
\frac{1}{(2\pi)^{3/2}}\left[\intl_{\Rnu^3}\vert \vec{k}\vert^{2(1+n)}\hat{\phi}(\vec{k})d\vec{k}\right]^{1/2}\mathcal{H}_X^{1/2}(\gamma),
\gamma\in\mathcal{D}_X,\label{eqn:VelocitySmoothness}
\end{equation}
provided that the integral
$\intl_{\Rnu^3}\vert \vec{k}\vert^{2(1+n)}\hat{\phi}(\vec{k})d\vec{k}$ is
finite and $\phi\in C^{n+4}(\Rnu^3,\Rnu^3)$.
\end{lemma}
\begin{proof}[Proof of Lemma \ref{lem:VelocitySmoothness}]
For smooth $\gamma$ Lemma has been proved in
\cite{[Bers+Gub_2007]} Lemma $3$. In general case, when
$\gamma\in\mathcal{D}_X$, it is enough to notice that both sides
of inequality \eqref{eqn:VelocitySmoothness} are locally Lipshitz
and therefore, continuous w.r.t. distance
$d(Y,\tilde{Y}):=\vert Y-\tilde{Y}\vert_D$, $Y\in\mathcal{D}_X$,
$\tilde{Y}\in\mathcal{D}_{\tilde{X}}$. Indeed, continuity of
$\mathcal{H}_X$ has been proven in Lemma
\ref{lem:EnergyProperties} and continuity of $\Vert \nabla^n
u^{\gamma}\Vert_{L^{\infty}}$ follows from Lemma
\ref{lem:ParametricRoughPath}.
\end{proof}
Now we are ready to prove Theorem \ref{thm:GlobalExistUniq}.
\begin{proof}[Proof of Theorem \ref{thm:GlobalExistUniq}]
According to Theorem \ref{thm:LocalExistUniq} there exists unique
local solution of problem
\eqref{eqn:RandomFilam-1}-\eqref{eqn:RandomFilam-3}. Then, we can
find $T^*>0$ such that there exists unique maximal local solution
$\gamma:[0,T^*)\to \mathcal{D}_{\gamma_0}$ and
\begin{equation}
\liml_{t\uparrow
T^*}\Vert \gamma(t)\Vert_{\mathcal{D}_{\gamma_0}}=\infty,\label{eqn:thmGEUaux-1}
\end{equation}
see e.g. \cite{[Cartan-1971]}. Notice that we will have
\begin{equation}
\frac{d\mathcal{H}_{\gamma_0}(\gamma(s))}{ds}=0,s\in[0,T^*).\label{eqn:EnergyMotionIntegral-1}
\end{equation}
Indeed, by Theorem \ref{thm:LocalExistUniq} for any $t_0\in
[0,T^*)$ there exists unique local solution $\tilde{\gamma}$ of
problem \eqref{eqn:RandomFilam-1}, \eqref{eqn:RandomFilam-3} with
initial condition $\gamma(t_0)$ on segment $[t_0,t_0+\eps_{t_0}]$
for some $\eps_{t_0}>0$. Therefore, $\gamma=\tilde{\gamma}$ on the
segment $[t_0,t_0+\eps_{t_0}]$. Hence,
$$
\frac{d\mathcal{H}_{\gamma_0}(\gamma(s))}{ds}=0,s\in[t_0,t_0+\eps_{t_0}],t_0\in[0,T^*),
$$
and identity \eqref{eqn:EnergyMotionIntegral-1} follows. We need
to show that $T^*=\infty$. Therefore, it is enough to prove
$$
\supl_{t\in[0,T^*)}\Vert \gamma(t)\Vert_{\mathcal{D}_{\gamma_0}}<\infty.
$$
Indeed, by contradiction with \eqref{eqn:thmGEUaux-1}, the result
will follow. In the rest of the proof we show such estimate. We
recall that
\begin{equation}
\gamma(t)=\gamma_0+\intl_0^tu^{\gamma(s)}(\gamma(s))ds.\label{eqn:RandomFilam-12}
\end{equation}
Firstly we have
\begin{align}
\vert \gamma(t)\vert_{L^{\infty}}&\leq\vert \gamma_0\vert_{L^{\infty}}+\intl_0^t\vert u^{\gamma(s)}\vert_{L^{\infty}}ds\nonumber\\
&\leq\vert \gamma_0\vert_{L^{\infty}}+C\intl_0^t\mathcal{H}_{\gamma_0}(\gamma(s))ds\nonumber\\
&\leq\vert \gamma_0\vert_{L^{\infty}}+C\mathcal{H}_{\gamma_0}(\gamma_0)t,t\in[0,T^*).\label{eqn:GEUEstim-1}
\end{align}
It follows from \eqref{eqn:RandomFilam-12} that
\begin{equation}
\gamma'(t)=\gamma_0'+\intl_0^t\nabla
u^{\gamma(s)}(\gamma(s))\gamma'(s)ds,t\in[0,T^*).\label{eqn:RandomFilam-12'}
\end{equation}
Therefore, by Lemmas \ref{lem:EnergyMotionIntegral} and
\ref{lem:VelocitySmoothness}
\begin{align}
\vert \gamma'(t)\vert_{L^{\infty}}&\leq\vert \gamma'_0\vert_{L^{\infty}}+\intl_0^t\vert \nabla
u^{\gamma(s)}\vert_{L^{\infty}}\vert \gamma'(s)\vert_{L^{\infty}}ds\nonumber\\
&\leq\vert \gamma'_0\vert_{L^{\infty}}+\intl_0^tC\mathcal{H}_{\gamma_0}^{1/2}(\gamma(s))\vert \gamma'(s)\vert_{L^{\infty}}ds\nonumber\\
&=\vert \gamma'_0\vert_{L^{\infty}}+\intl_0^tC\mathcal{H}_{\gamma_0}^{1/2}(\gamma_0)\vert \gamma'(s)\vert_{L^{\infty}}ds,t\in[0,T^*).
\end{align}
Then by Gronwall inequality we infer our second estimate
\begin{equation}
\vert \gamma'(t)\vert_{L^{\infty}}\leq\vert \gamma'_0\vert_{L^{\infty}}e^{C\mathcal{H}_{\gamma_0}^{1/2}(\gamma_0)t},t\in[0,T^*).\label{eqn:GEUEstim-2}
\end{equation}
We will need one auxiliary estimate. We have
\begin{align}
\vert \gamma(t)\vert_{C^{\nu}}&\leq\vert \gamma_0\vert_{C^{\nu}}+\intl_0^t\vert
u^{\gamma(s)}(\gamma(s))\vert_{C^{\nu}}\vert ds\nonumber\\
&\leq\vert \gamma_0\vert_{C^{\nu}}+\intl_0^t\vert \nabla
u^{\gamma(s)}\vert_{L^{\infty}}\vert \gamma(s)\vert_{C^{\nu}}ds\nonumber\\
&\leq
\vert \gamma_0\vert_{C^{\nu}}+\intl_0^tC\mathcal{H}_{\gamma_0}^{1/2}(\gamma(s))\vert \gamma(s)\vert_{C^{\nu}}ds\nonumber\\
&=\vert \gamma_0\vert_{C^{\nu}}+\intl_0^tC\mathcal{H}_{\gamma_0}^{1/2}(\gamma_0)\vert \gamma(s)\vert_{C^{\nu}}ds,t\in[0,T^*).
\end{align}
Thus, by Gronwall inequality we get
\begin{equation}
\vert \gamma(t)\vert_{C^{\nu}}\leq\vert \gamma_0\vert_{C^{\nu}}e^{C\mathcal{H}_{\gamma_0}^{1/2}(\gamma_0)t},t\in[0,T^*).\label{eqn:GEUEstim-aux}
\end{equation}
Now we can estimate $C^{\nu}$ norm of $\gamma'$. We have
\begin{align}
\vert \gamma'(t)\vert_{C^{\nu}}&\leq\vert \gamma'_0\vert_{C^{\nu}}+\intl_0^t\vert \nabla
u^{\gamma(s)}(\gamma(s))\gamma'(s)\vert_{C^{\nu}}ds\nonumber\\
&\leq \vert \gamma'_0\vert_{C^{\nu}}+\intl_0^t(\vert \nabla
u^{\gamma(s)}\vert_{L^{\infty}}\vert \gamma'(s)\vert_{C^{\nu}}+\vert \gamma'(s)\vert_{L^{\infty}}\vert \nabla
u^{\gamma(s)}(\gamma(s))\vert_{C^{\nu}})ds\nonumber\\
&\leq\vert \gamma'_0\vert_{C^{\nu}}+\intl_0^t(\vert \nabla
u^{\gamma(s)}\vert_{L^{\infty}}\vert \gamma'(s)\vert_{C^{\nu}}+\vert \gamma'(s)\vert_{L^{\infty}}\vert \nabla^2
u^{\gamma(s)}\vert_{L^{\infty}}\vert \gamma(s)\vert_{C^{\nu}})ds\nonumber\\
&\leq\vert \gamma'_0\vert_{C^{\nu}}\nonumber\\
&+\intl_0^t(C\mathcal{H}_{\gamma_0}^{1/2}(\gamma_0)(\vert \gamma'(s)\vert_{C^{\nu}}+\vert \gamma'_0\vert_{L^{\infty}}\vert \gamma_0\vert_{C^{\nu}}e^{C\mathcal{H}_{\gamma_0}(\gamma_0)s}))ds,t\in[0,T^*),
\end{align}
where last inequality follows from Lemmas
\ref{lem:EnergyMotionIntegral} and \ref{lem:VelocitySmoothness}.
Then by Gronwall inequality we get third estimate
\begin{equation}
\vert \gamma'(t)\vert_{C^{\nu}}\leq(\vert \gamma'_0\vert_{C^{\nu}}+\vert \gamma'_0\vert_{L^{\infty}}\vert \gamma_0\vert_{C^{\nu}})e^{C\mathcal{H}_{\gamma_0}(\gamma_0)t},t\in[0,T^*).\label{eqn:GEUEstim-3}
\end{equation}
%and by inequality \eqref{eqn:GEUEstim-aux} we get third estimate
%\begin{equation}
%\vert \gamma'(t)\vert_{C^{\nu}}\leq\vert \gamma'_0\vert_{C^{\nu}}e^{C\mathcal{H}_X^{1/2}(\gamma_0)te^{C\mathcal{H}_X^{1/2}(\gamma_0)t}},t\in[0,T^*).\label{eqn:GEUEstim-3}
%\end{equation}
It remains to find an estimate for $\vert R^{\gamma(t)}\vert_{2\nu}$. We
have
\begin{equation}
R^{\gamma(t)}=R^{\gamma_0}+\intl_0^tR^{u^{\gamma(s)}(\gamma(s))}ds,t\in[0,T^*).
\end{equation}
By identity \eqref{eqn:RegularMap-2} we have for $s\in[0,T^*)$
\begin{align}
R^{u^{\gamma(s)}(\gamma(s))}(\xi,\eta)&=\nabla
u^{\gamma(s)}(\gamma(s,\xi))R^{\gamma(s)}(\xi,\eta)+
\suml_k(\gamma^k(s,\eta)-\gamma^k(s,\xi))\times\nonumber\\
&\intl_0^1\left[\frac{\partial u^{\gamma(s)}}{\partial
x_k}(\gamma(s,\xi)+r(\gamma(s,\eta)-\gamma(s,\xi)))-\frac{\partial
u^{\gamma(s)}}{\partial
x_k}(\gamma(s,\xi))\right]dr.\label{eqn:GEUEstim-aux3}
\end{align}
Therefore,
\begin{equation}
\vert R^{u^{\gamma(s)}(\gamma(s))}\vert_{\tilde{C}^{2\nu}}\leq \vert \nabla
u^{\gamma(s)}\vert_{L^{\infty}}\vert R^{\gamma(s)}\vert_{\tilde{C}^{2\nu}}+\frac{1}{2}\vert \gamma(s)\vert_{C^{\nu}}^2\vert \nabla^2
u^{\gamma(s)}\vert_{L^{\infty}},s\in[0,T^*).\label{eqn:GEUEstim-aux4}
\end{equation}
Thus, by inequalities \eqref{eqn:GEUEstim-aux4} and
\eqref{eqn:GEUEstim-aux} we have for $t\in[0,T^*)$
\begin{align}
\vert R^{\gamma(t)}\vert_{\tilde{C}^{2\nu}}&\leq\vert R^{\gamma_0}\vert_{\tilde{C}^{2\nu}}+\intl_0^t(\vert \nabla
u^{\gamma(s)}\vert_{L^{\infty}}\vert R^{\gamma(s)}\vert_{\tilde{C}^{2\nu}}+\frac{1}{2}\vert \gamma(s)\vert_{C^{\nu}}^2\vert \nabla^2
u^{\gamma(s)}\vert_{L^{\infty}})ds\nonumber\\
&\leq\vert R^{\gamma_0}\vert_{\tilde{C}^{2\nu}}+\intl_0^t(\vert \nabla
u^{\gamma(s)}\vert_{L^{\infty}}\vert R^{\gamma(s)}\vert_{\tilde{C}^{2\nu}}+\vert \gamma_0\vert_{C^{\nu}}e^{C\mathcal{H}_{\gamma_0}^{1/2}(\gamma_0)t}\vert \nabla^2
u^{\gamma(s)}\vert_{L^{\infty}})ds\nonumber\\
&\leq\vert R^{\gamma_0}\vert_{\tilde{C}^{2\nu}}+C(\vert \gamma_0\vert_{C^{\nu}},\mathcal{H}_{\gamma_0}(\gamma_0))e^{C\mathcal{H}_{\gamma_0}^{1/2}(\gamma_0)t}\nonumber\\
&+\intl_0^tC\mathcal{H}_{\gamma_0}^{1/2}(\gamma_0)\vert R^{\gamma(s)}\vert_{\tilde{C}^{2\nu}}ds,\label{eqn:GEUEstim-aux5}
\end{align}
where in the last inequality we used Lemmas
\ref{lem:EnergyMotionIntegral} and \ref{lem:VelocitySmoothness}.
As the result, by Gronwall Lemma we get
\begin{equation}
\vert R^{\gamma(t)}\vert_{\tilde{C}^{2\nu}}\leq(\vert R^{\gamma_0}\vert_{\tilde{C}^{2\nu}}+C(\vert \gamma_0\vert_{C^{\nu}},\mathcal{H}_{\gamma_0}(\gamma_0))e^{C\mathcal{H}_{\gamma_0}^{1/2}(\gamma_0)t}))
e^{C\mathcal{H}_{\gamma_0}^{1/2}(\gamma_0)t},t\in[0,T^*),\label{eqn:GEUEstim-4}
\end{equation}
and combining estimates \eqref{eqn:GEUEstim-1},
\eqref{eqn:GEUEstim-2}, \eqref{eqn:GEUEstim-3}, and
\eqref{eqn:GEUEstim-4} we prove following a'priori estimate
\begin{equation}
\vert \gamma(t)\vert_{\mathcal{D}_{\gamma_0}}\leq
K(1+\mathcal{H}_{\gamma_0}(\gamma_0))(1+\vert \gamma_0\vert_{\mathcal{D}_{\gamma_0}})\vert \gamma_0\vert_{\mathcal{D}_{\gamma_0}}e^{C\mathcal{H}_{\gamma_0}(\gamma_0)t},t\in[0,T^*),\label{eqn:GEUEstim-5}
\end{equation}
and the result follows.
\end{proof}

\section{Future directions of research}

It would be interesting to consider problem
\eqref{eqn:RandomFilam-1}-\eqref{eqn:RandomFilam-3} with added
white noise i.e. to consider problem
\begin{eqnarray}
d\gamma(t) &=& u^{\gamma(t)}(\gamma(t))dt+\sqrt{2\nu}dw_t,\nu>0,t\in [0,T] \label{eqn:FutureRandomFilam-1}\\
\gamma(0) &=& \gamma_0,\label{eqn:FutureRandomFilam-2}
\end{eqnarray}
where $\gamma_0$ is a geometric $\nu$-rough path, vector field of velocity $u^Y$ is given by
\eqref{eqn:RandomFilam-2} and $w_t$ is $\mathcal{D}_{\gamma_0}$--valued
Wiener process. This model would correspond to Navier-Stokes
equations rather than Euler equations. There are two possible
mathematical frameworks for the model.

First one is to make change of variables
$$
\alpha(t)=\gamma(t)-\sqrt{2\nu}w_t, t\in[0,T].
$$
Then we can fix $\{w_t\}_{t\geq 0}$ and system
\eqref{eqn:FutureRandomFilam-1}-\eqref{eqn:FutureRandomFilam-2} is
reformulated as follows
\begin{eqnarray}
\frac{d\alpha}{dt} &=& u^{\alpha(t)+\sqrt{2\nu}w_t}(\alpha(t)+\sqrt{2\nu}w_t),\nu>0,t\in [0,T] \label{eqn:FutureRandomFilam-1a}\\
\alpha(0) &=& \gamma_0.\label{eqn:FutureRandomFilam-2a}
\end{eqnarray}
Now the problem
\eqref{eqn:FutureRandomFilam-1a}-\eqref{eqn:FutureRandomFilam-2a}
is ordinary differential equation (ODE) with random coefficients
in $\mathcal{D}_{\gamma_0}$ and it can be studied by methods of the theory
of random dynamical systems, see \cite{[Arnold-98]} and
\cite{[DuanLuSchmalfus-2003]}. This approach works only in the
case of additive noise.

Second approach is to consider problem
\eqref{eqn:FutureRandomFilam-1}-\eqref{eqn:FutureRandomFilam-2} as
SDE in Banach space $\mathcal{D}_{\gamma_0}$. Then, we can consider more
general system with multiplicative noise:
\begin{eqnarray}
d\gamma(t) &=& u^{\gamma(t)}(\gamma(t))dt+\sqrt{2\nu}G(\gamma)dw_t,\nu>0,t\in [0,T] \label{eqn:FutureRandomFilam-1b}\\
\gamma(0) &=& \gamma_0.\label{eqn:FutureRandomFilam-2b}
\end{eqnarray}
The problem which appear here is to define Stochastic integral in
the Banach space $\mathcal{D}_{\gamma_0}$. Stochastic calculus in M-type 2
Banach spaces developed in works
\cite{[Dettweiler-85]}-\cite{[Dettweiler-91]}, \cite{[Brz-95]},
\cite{[Brz-03]} does not work in this situation. It seems that it
is necessary to try to alter definition of $\mathcal{D}_{\gamma_0}$ to be
able to apply the theory.

Other possible direction of research is the theory of connections
on infinite dimensional manifolds, see \cite{[Gawedzki-1988]},
\cite{[Brylinski-1992]}, \cite{[Leandre-1999]}. In
\cite{[Brz-Leandre-2000]} the authours claimed, see p. 251
therein, that it is possible to define the topological space of
Gaw{\c e}dzki's \cite{[Gawedzki-1988]} line bundle over the set of
rough loops in the sense of Lyons \cite{[Lyons-1998]}. Since the
trajectories of the Brownian loop are almost surely rough paths,
this allows us to define the topological space of Gaw{\c e}dzki's
line bundle over the Brownian bridge, because it is possible to
define the integral of a one-form over a rough path. It would be
interesting to write down a complete proof of this claim. The
theory presented in this article could be seen as a first step in
realizing such a programme.

\begin{thebibliography}{15}

\bibitem{[Arnold-98]} L. Arnold,
\newblock\textit{ Random dynamical systems,}
\newblock{Springer-Verlag, Berlin Heidelberg, New York, 1998}


\bibitem{[BKM-1984]} J. T. Beale, T. Kato, and A. Majda,
\textit{ Remarks on the breakdown of smooth solutions for the 3-D
Euler equations}, Comm. Math. Phys. \textbf{94}, no. 1, 61--66
(1984)

\bibitem{[Bell+Markus_1992]} J. Bell and D. Markus,
\textit{ Vorticity intensification and the transition to
turbulence in the three-dimensional Euler equation},Comm. Math.
Phys. \textbf{147}, no. 2, 371--394 (1992)


\bibitem{[Bers+Gub_2007]} L. C. Berselli, M. Gubinelli,
\textit{ On the global evolution of vortex filaments, blobs, and
small loops in 3D ideal flows}, Commun. Math. Phys., \textbf{269}, 693--713 (2007)

\bibitem{[Bers+Bess_2002]} L. C. Berselli, H. Bessaih,
\textit{ Some results for the line vortex equation}, Nonlinearity,
\textbf{15}, no. 6, 1729--1746 (2002)

\bibitem{[Bess+Gub_2005]} H. Bessaih, M. Gubinelli, F.
Russo, \textit{ The evolution of a random vortex filament}, Ann.
Probab. \textbf{33}, no. 5, 1825--1855 (2005)

\bibitem{[Gubinelli-2004]} M. Gubinelli,
\textit{ Controlling rough paths,} Journal of Functional Analysis,
Vol.\textbf{216},86--140 (2004)

\bibitem{[Brz-95]} Z. Brze\'zniak,
\newblock{Stochastic partial differential equations in M-type $2$ Banach
spaces, \textit{ Potential Anal.} 4, no. 1, 1--45 (1995)}

\bibitem{[Brz-03]} Z. Brze\'{z}niak,
\newblock{Some remarks on It\^{o} and Stratonovich integration in $2$-smooth Banach spaces,}
\newblock{\textit{ Probabilistic methods in fluids}, proceedings of the Swansea 2002 workshop, 14-19, April 2002, Wales,
UK, World Scientific, 48-69 (2003)}

\bibitem{[Brz-Leandre-2000]} Z. Brze\'{z}niak, R. L\'{e}andr\'{e},
\textit{ Horizontal Lift of an Infinite Dimensional Diffusion},
Potential Anal. \textbf{12}, 249--280 (2000)

\bibitem{[Brylinski-1992]} J. L. Brylinski,
\newblock{\textit{Loop spaces, characteristic classes and geometric
quantization}, Progress. in Math. \textbf{107}, Birkhauser, Basel,
1992.}

\bibitem{[Cartan-1971]} H. Cartan,
\newblock\textit{ Differential Calculus},
\newblock{Hermann, Paris, 1971}

\bibitem{[Dettweiler-85]} E. Dettweiler,
\textit{ Stochastic integration of Banach space valued functions},
in \textit{ Stochastic Space-Time Models and Limit Theorem}, pp.
33-79, D.Reidel Publ. Co., 1985.

\bibitem{[Dettweiler-91]} E. Dettweiler,
\textit{ Stochastic Integration relative to Brownian Motion on a
general Banach space}, Boga-- Tr. J. of Mathematics \textbf{15},
6--44 (1991)

\bibitem{[DuanLuSchmalfus-2003]} J. Duan, K. Lu, B. Schmalfuss,
\textit{ Invariant manifolds for stochastic partial differential
equations}, The Annals of Probability \textbf{31}, no. 4,
2109--2135 (2003)

\bibitem{[Friz-2005]} P. Friz,
\newblock{\textit{ Continuity of the It\^o-map for Holder rough paths with
applications to the support theorem in Holder norm}, Probability
and partial differential equations in modern applied mathematics,
117--135, IMA Vol. Math. Appl., \textbf{140}, Springer, New York,
2005.}

\bibitem{Gallavotti} G. Gallavotti,
\newblock{\sc Foundations of fluid dynamics. Translated from the Italian.,}
\newblock{Texts and Monographs in Physics. Springer-Verlag, Berlin, 2002.}

\bibitem{[Gawedzki-1988]} K. Gaw{\c e}dzki,
\newblock{\textit{ Topological actions in two-dimensional quantum field
theory}, In: Gooft G�t., Jaffe A., Mack G., Mitter P. K., Stora R.
(eds.),Nonperturbative Quantum Field Theory, NATO series
\textbf{185}, Plenum Press, 101�-142, 1988}

\bibitem{Helm1858} H. Helmholtz,
\textit{\"{U}ber integrale der hydrodynamischen gleichungen, welche den Wirbelbewegungen
entsprechen}, Crelle J. \textbf{55}, 25 (1858)

\bibitem{Kelvin1883} Lord Kelvin,\textit{On vortex motion}, Trans. R. Soc. Edin. \textbf{25}, 217--260 (1883)

\bibitem{[Leandre-1999]} R. L\'{e}andr\'{e},
\textit{ String Structure over the Brownian Bridge}, J. Math.
Phys. \textbf{40}, no. 1, 454--479 (1999)

\bibitem{[Lyons-1998]} T. J. Lyons,
\textit{ Differential equations driven by rough signals}, Rev.
Mat. Iberoamericana, \textbf{14}, no. 2, 215--310 (1998)

\bibitem{[LyonsQian-2002]} T. J. Lyons, Z. Qian,
\newblock{\sc System Control and Rough Paths,}
\newblock{Oxford University Press, Oxford, 2002.}

\bibitem{[Rosenhead-1930]} L. Rosenhead,
\textit{ The spread of vorticity in the wake behind a cylinder},
Proc. Royal Soc. \textbf{127}, 590--612 (1930)


\bibitem{[VincMeneguzzi-1991]} A. Vincent, M. Meneguzzi,
\textit{ The spatial structure and the statistical properties of
homogeneous turbulence}, J. Fluid. Mech. \textbf{225}, no.1, 1--25
(1991)

\bibitem{[Young-1936]} L. C. Young,
\textit{ An inequality of H\"older type, connected with Stieljes
integration}, Acta Math. \textbf{67}, 251--282 (1936)

\end{thebibliography}
\end{document}
